% !TeX root = ../main.tex

\chapter{Introduction}
\label{ch:introduction}

\section{Background and Motivation}
\label{sec:intro-background}

Metagenomics, the culture-independent study of microbial communities through direct sequencing of environmental DNA, has transformed our understanding of microbial ecology, human health, and environmental science~\cite{Quince2017}. 
Modern metagenomic sequencing platforms, particularly Illumina short-read technologies, generate hundreds of millions of reads per sample, each typically 100--300 base pairs (bp) in length. 
A fundamental computational challenge in metagenomic analysis is \emph{taxonomic classification}: assigning each sequencing read to the correct taxon in the microbial taxonomy, such as genus or species~\cite{Ye2019}.

Accurate taxonomic classification underpins virtually all downstream metagenomic analyses, including community composition profiling, pathogen detection, antibiotic resistance gene tracking, and comparative microbiome studies~\cite{Breitwieser2019}. 
However, the classification task is inherently difficult due to several factors:
\begin{itemize}
    \item \textbf{Short read length}: Individual reads of 150~bp carry limited phylogenetic signal, especially when distinguishing closely related taxa.
    \item \textbf{High taxonomic diversity}: A typical metagenomic sample may contain reads from hundreds to thousands of species spanning dozens of genera.
    \item \textbf{Class imbalance}: Community compositions follow highly skewed abundance distributions, with a few dominant taxa and many rare ones.
    \item \textbf{Sequence similarity}: Closely related taxa share substantial genomic homology, making fine-grained discrimination (e.g., species-level) particularly challenging.
\end{itemize}

Traditional approaches to metagenomic classification fall broadly into two categories. 
\emph{Alignment-based methods} such as BLAST~\cite{Altschul1990} and DIAMOND~\cite{Buchfink2015} align reads against reference databases, achieving high precision but at considerable computational cost. 
\emph{$k$-mer-based methods} such as Kraken2~\cite{Wood2019} and Centrifuge~\cite{Kim2016} index reference genomes using $k$-mer hash tables, enabling ultra-fast classification at the expense of sensitivity. 
More recently, deep learning approaches have emerged, including MetaTransformer~\cite{Wichmann2023}, which applies a Transformer architecture to learned $k$-mer embeddings for metagenomic read classification.

Concurrently, the field of genomics has witnessed the development of \emph{genomic foundation models} (GFMs)---large-scale pre-trained models that learn general-purpose representations of DNA sequences from billions of nucleotides~\cite{DallaTorre2023}. 
Models such as the Nucleotide Transformer~\cite{DallaTorre2023}, DNABERT~\cite{Ji2021}, DNABERT-2~\cite{Zhou2024dnabert2}, and Evo~\cite{Nguyen2024} have demonstrated strong performance across diverse genomic tasks, including promoter prediction, splice site identification, and variant effect prediction. 
These models offer a compelling alternative to task-specific feature engineering: by pre-training on vast genomic corpora, they capture rich sequence semantics that can be transferred to downstream tasks through fine-tuning.

While GFMs have been evaluated extensively on promoter, enhancer, splice-site, TF-binding, and other molecular phenotype prediction tasks, their behavior in high-cardinality metagenomic short-read taxonomic classification remains poorly characterized.
A na\"ive approach---extracting mean-pooled embeddings from a frozen pre-trained model and training a linear classifier---discards the positional and local information encoded in individual token representations.
A separate but equally important question concerns \emph{tokenization}: most existing GFMs (including NT-v2) use non-overlapping or short-$k$ tokenization schemes inherited from masked-language-model pre-training, yet $k$-mer-based metagenomic classifiers commonly rely on overlapping long $k$-mers (e.g., $k = 12$--13) for taxonomic specificity.
It is therefore unclear whether large-scale GFM pre-training can compensate for tokenization choices that are mismatched to the downstream taxonomic task.

This thesis investigates two complementary questions: (i)~whether a \emph{token-level} approach, preserving the full sequence of token embeddings and employing parameter-efficient fine-tuning, can substantially improve metagenomic taxonomic classification over frozen-embedding baselines; and (ii)~whether GFM pre-training is sufficient to compensate for the absence of task-specific $k$-mer tokenization, or whether tokenization design imposes a structural ceiling that pre-training alone cannot overcome.

\section{Problem Statement}
\label{sec:intro-problem}

Given a set of metagenomic short reads $\{r_1, r_2, \ldots, r_N\}$ generated by sequencing a microbial community, the taxonomic classification problem is to assign each read $r_i$ to its correct taxon $y_i \in \{1, 2, \ldots, K\}$, where $K$ is the number of classes at a given taxonomic level (e.g., genus or species).

This thesis formulates the problem as a supervised multi-class classification task. We assume access to a reference genome database from which simulated reads with known taxonomic labels can be generated for training. The challenge is to learn a classifier $f_\theta: \mathcal{R} \rightarrow \{1, \ldots, K\}$ that generalizes well from simulated training data to accurately classify novel reads.

The specific objectives of this thesis are:
\begin{enumerate}
    \item \textbf{Token-level representation}: Investigate whether preserving the full sequence of token embeddings from a genomic foundation model, rather than using mean-pooled summaries, improves taxonomic classification accuracy.
    \item \textbf{Parameter-efficient fine-tuning}: Evaluate the effectiveness of LoRA (Low-Rank Adaptation)~\cite{Hu2022} for adapting a 498M-parameter pre-trained model to the taxonomic classification task with minimal trainable parameters.
    \item \textbf{Data scaling}: Quantify the relationship between training data volume and classification performance within a fixed GFM-based architecture, and determine the data requirements for achieving competitive accuracy.
    \item \textbf{Pre-training vs.\ tokenization}: Through controlled cross-architecture ablations holding training data and capacity comparable, isolate the contributions of (a)~large-scale genomic pre-training and (b)~$k$-mer tokenization design to taxonomic-classification performance, and identify the dominant bottleneck of the proposed approach.
    \item \textbf{Multi-level and sample-level evaluation}: Assess performance across multiple metrics (micro accuracy, macro F1, top-$k$ accuracy) at both genus and species taxonomic levels, and characterize how per-read accuracy translates to sample-level utility (relative-abundance estimation, presence/absence detection).
\end{enumerate}

\section{Contributions}
\label{sec:intro-contributions}

The main contributions of this thesis are as follows:

\begin{enumerate}
    \item \textbf{Systematic evaluation of GFM fine-tuning for metagenomic classification}: We propose a token-level classification pipeline that feeds the complete per-token output of NT-v2 into a learnable attention-pooling head and fine-tunes the backbone with LoRA, covering three GFMs (NT-v2, DNABERT, DNABERT-2) across frozen-embedding baselines, parameter-efficient fine-tuning, data scaling, RC test-time augmentation, and sample-level utility. This is, to our knowledge, the first end-to-end empirical study of GFM fine-tuning for short-read taxonomic classification at the 120-genus and 1,535-species scale.

    \item \textbf{Data scaling within a fixed GFM-based architecture}: We conduct systematic data-scaling experiments from 500K to 250M reads within the fixed NT-v2 6-mer framework. The data reveal a log-saturation pattern: $500\text{K}\to 5\text{M}$ yields $+7.76$~pp genus accuracy, $5\text{M}\to 50\text{M}$ adds $+4.02$~pp, yet a further $5\times$ increase to 250M contributes only $+0.22$~pp---empirically refuting log-linear scaling projections. This saturation is tokenization-bound rather than data-bound: an overlapping 13-mer model trained on the same data range continues scaling from 87.5\% to 98.7\%, confirming that NT-v2's 6-mer tokenizer is the binding constraint.

    \item \textbf{KmerFormer, and what pre-training is actually worth}: We build KmerFormer, a from-scratch $k$-mer transformer in which the tokenizer, the stride, the vocabulary and the encoder depth are free under one training recipe, so that arms differing in one of them stay comparable. Using it we sweep depth at NT-v2's own tokenizer and find that depth alone is worth $+15.60$~pp (53.92\% to 69.52\% from 1 to 29 layers). This corrects the attribution an earlier single-layer ablation invited: at matched depth and 50M labelled reads, pre-training's net contribution is approximately zero, a 6.45M-parameter from-scratch model exceeding a 500M pre-trained one by 2.44~pp. What pre-training does buy is \emph{data efficiency}---$+13.4$, $+13.1$ and $+8.9$~pp at 0.5M, 1M and 5M reads, crossing over by 50M---and the crossover reproduces on an independent 309-genus soil catalogue.

    \item \textbf{A long-$k$ vocabulary in our own codebase, and the cost of approximating it}: We train 13-mer arms under the same recipe as the 6-mer ones, so the tokenization effect is measured without crossing a codebase boundary. An exact canonical vocabulary reaches 91.15\% genus Top-1, and a hashed vocabulary that maps each 13-mer into $2^{22}$ buckets reaches 86.80\% while replacing an 8~GiB embedding table with a 2~GiB one---and, unlike an enumerated table, trains under data parallelism at all. Hashing costs 6.37~pp at matched width and depth. Depth reverses sign at this tokenization: one layer beats sixteen, in both vocabularies.

    \item \textbf{Tokenization is the dominant design factor}: Holding MetaTransformer's architecture fixed while varying $k \in \{6, 12, 13\}$ and stride exposes two compounding effects: overlapping stride ($s{=}1$) versus non-overlapping ($s{=}k$) accounts for up to $+60$~pp at $k{=}13$, and longer $k$ improves accuracy monotonically, reaching $+40.42$~pp at 13-mer. Measured against the $+15.60$~pp that depth is worth inside our own codebase, tokenization is the larger factor by roughly $2.6\times$---and both estimates are obtained without crossing a codebase boundary. A third factor is \emph{not} what an earlier version of this work reported: the NT-v2-against-MetaTransformer difference is not a single-factor measurement of pre-training, and that decomposition is withdrawn in Section~\ref{sec:exp-tokenization}.

    \item \textbf{Closed-set accuracy measures recognition, not generalization}: We evaluate every arm on genomes absent from the reference, stratified by ANI into strain, species and far rungs. The ranking inverts along that gradient: a 13-mer arm leads by 32 points on a new strain of a catalogued species and trails by 11.3 at the far rung, where a 6-mer model whose 6.45M parameters are a 333rd of the 13-mer arm's 2.15\,B reaches 22.37\% against 11.07\%. Retention separates the two vocabularies into disjoint bands (49--54\% against 13--22\%) that no difference of depth, capacity or pre-training crosses, and an embedding-space measurement taken without any classifier reproduces the same split. The 98.7\% attainable under complete reference coverage is therefore an upper bound obtained by recognising catalogued genomes.

    \item \textbf{Reverse-complement test-time augmentation}: We evaluate RC TTA---averaging logits from forward and reverse-complement strand passes at inference---across all experimental settings and architectures. RC TTA delivers a zero-training-cost improvement of $+0.1$ to $+1.5$~pp for pre-trained backbones and approximately zero for from-scratch models at \emph{every} depth. A pair of 13-mer arms locates the mechanism in the vocabulary rather than in pre-training: a canonical vocabulary, in which a $k$-mer and its reverse complement share an entry, gains $+0.01$~pp because the augmentation is then arithmetically a no-op, while a non-canonical hashed vocabulary gains $+6.71$~pp.

    \item \textbf{Hierarchical species classification and sample-level utility}: We extend the evaluation framework to species-level classification (1,535 classes) via a hierarchical genus-routing pipeline, and to sample-level abundance estimation. Per-genus LoRA specialists under oracle-genus routing reach 29.5\% species Top-1, exceeding the 27.4\% flat-classifier oracle and validating hierarchical specialisation; the genus-router accuracy (66\%--67\%) is identified as the binding bottleneck for end-to-end routing. At the sample level, 67\% per-read accuracy on species-balanced synthetic samples translates to Pearson $r = 0.993$ for relative abundance and 100\% detection sensitivity at $\geq$1\% abundance, though generalization to real and compositionally variable communities remains future work.
\end{enumerate}

\section{Thesis Organization}
\label{sec:intro-organization}

The remainder of this thesis is organized as follows:

\begin{itemize}
    \item \textbf{Chapter~\ref{ch:background}} reviews the background and related work, covering metagenomic classification methods, genomic foundation models, parameter-efficient fine-tuning, and attention mechanisms.
    
    \item \textbf{Chapter~\ref{ch:methodology}} presents the proposed methodology in detail, including the model architecture, training strategy, data pipeline, and evaluation framework.
    
    \item \textbf{Chapter~\ref{ch:experiments}} describes the experimental setup and presents comprehensive results, including ablation studies, data scaling analysis, species-level evaluation, and computational resource assessment.
    
    \item \textbf{Chapter~\ref{ch:conclusion}} summarizes the findings, discusses limitations, and outlines directions for future work.
\end{itemize}
